\documentclass[11pt,a4paper]{article}

\usepackage[utf8]{inputenc}
\usepackage[T1]{fontenc}
\usepackage[brazil]{babel}
\usepackage[a4paper,margin=2.4cm]{geometry}
\usepackage{listings}
\usepackage{xcolor}
\usepackage{enumitem}
\usepackage[hidelinks]{hyperref}

\hypersetup{
  colorlinks=true,
  linkcolor=black,
  urlcolor=blue!55!black,
  citecolor=black,
}

\definecolor{codebg}{gray}{0.96}
\lstset{
  basicstyle=\ttfamily\footnotesize,
  backgroundcolor=\color{codebg},
  breaklines=true,
  columns=fullflexible,
  keepspaces=true,
  frame=single,
  rulecolor=\color{gray!40},
  framesep=5pt,
  xleftmargin=6pt,
  xrightmargin=6pt,
  showstringspaces=false,
  aboveskip=8pt,
  belowskip=8pt,
}

\setlist{itemsep=2pt,topsep=3pt}

\title{Infraestrutura do Surfer e do Surfer-Aurora}
\author{NIPSCERN}
\date{\today}

\begin{document}

\maketitle

\begin{abstract}
\noindent
Este documento descreve, de forma geral, a infraestrutura do visualizador de
formas de onda Surfer e do nosso fork Surfer-Aurora. Cobre o que o Surfer é e
com que tecnologia foi feito, a hierarquia de pastas do repositório no GitLab,
onde fica cada componente (interface, lógica, dados, formatadores, assets), os
links relevantes, como o nosso executável próprio é de fato gerado e o que
precisamos fazer para gerá-lo, e como a AURORA usa o \texttt{surfer-aurora.exe}.
\end{abstract}

\tableofcontents
\bigskip

\section{Visão geral do Surfer}

O Surfer é um visualizador de formas de onda (waveform viewer) com foco em uma
interface rápida e utilizável e em extensibilidade. Ele abre traços de sinais
digitais produzidos por simulação e permite inspecioná-los no tempo, com
curadoria de sinais, cores, formatos e grupos.

Pontos principais:

\begin{itemize}
  \item Escrito em Rust. A interface usa o \texttt{egui} (uma biblioteca de GUI
        em modo imediato) através do \texttt{eframe}, e o desenho é feito na GPU
        via \texttt{wgpu}.
  \item Lê os formatos VCD, FST e GHW. A leitura dos traços é feita pela
        biblioteca \texttt{wellen}, que usa uma estrutura de dados inspirada no
        FST para reduzir tempo de carga e uso de RAM em waveforms grandes.
  \item Roda nativo em Windows, macOS e Linux, roda no navegador via WebAssembly
        (demo em \url{https://app.surfer-project.org}) e existe como extensão do
        VS Code.
  \item É software livre sob a licença EUPL-1.2.
  \item É extensível: além dos formatadores de valor embutidos, aceita plugins
        (translators) e possui um protocolo de controle externo, o WCP (Waveform
        Control Protocol), além de um modo cliente-servidor (\texttt{surver}).
\end{itemize}

\section{Links}

\subsection{Projeto original (upstream)}
\begin{itemize}
  \item Repositório: \url{https://gitlab.com/surfer-project/surfer}
  \item Site e blog: \url{https://surfer-project.org} e
        \url{https://blog.surfer-project.org}
  \item Demonstração no navegador: \url{https://app.surfer-project.org}
  \item Documentação de usuário (livro):
        \url{https://docs.surfer-project.org/book/}
  \item Documentação do código-fonte:
        \url{https://docs.surfer-project.org/surfer/}
  \item Extensão do VS Code:
        \url{https://marketplace.visualstudio.com/items?itemName=surfer-project.surfer}
  \item Biblioteca de parsing \texttt{wellen}:
        \url{https://crates.io/crates/wellen}
  \item Canal Matrix da comunidade:
        \url{https://matrix.to/#/#surfer-project:matrix.org}
\end{itemize}

\subsection{Nosso fork (NIPSCERN)}
\begin{itemize}
  \item Grupo no GitLab: \url{https://gitlab.com/nips-cern}
  \item Fork Surfer-Aurora: \url{https://gitlab.com/nips-cern/surfer-aurora}
  \item Registro de pacotes (onde publicamos o binário):
        \url{https://gitlab.com/nips-cern/surfer-aurora/-/packages}
\end{itemize}

\section{Estrutura do repositório}

O repositório é um workspace do Cargo (o gerenciador de pacotes do Rust): vários
crates relacionados dentro de um mesmo repositório, coordenados pelo
\texttt{Cargo.toml} da raiz. A hierarquia de pastas, com os nomes reais, é a
seguinte.

\subsection{Raiz}

\begin{lstlisting}
surfer-aurora/
  Cargo.toml                 workspace: lista os crates e as dependencias
  Cargo.lock                 versoes travadas (build reproduzivel)
  README.md  CHANGELOG.md
  LICENSE-EUPL-1.2.txt        licenca do projeto (EUPL-1.2)
  default_config.toml         configuracao padrao do Surfer
  default_theme.toml          tema padrao
  book.toml                   configuracao do livro (mdBook) das docs
  surfer.sh                   script auxiliar (rodar o Surfer a partir do WSL)

  libsurfer/                  o NUCLEO: toda a logica e a interface
  surfer/                     o binario desktop/web (ponto de entrada da GUI)
  surver/                     o servidor headless (modo cliente-servidor)
  surfer-translation-types/   tipos compartilhados dos formatadores/plugins
  surfer-wcp/                 o Waveform Control Protocol
  surfer-vscode/             a extensao do VS Code
  wasm_example_translator/    exemplo de plugin (translator em WASM)

  translator-docs/            documentacao dos formatadores
  themes/                     temas prontos (.toml)
  docs/                       fonte da documentacao de usuario
  examples/                   waveforms de exemplo (.vcd/.fst)
  snapshots/                  imagens de referencia dos testes visuais
  tools/                      scripts utilitarios (Python)
  misc/                       arquivos diversos
  target/                     saida de build (gerada; nao versionada)
\end{lstlisting}

\subsection{Crates do workspace}

O workspace declara os membros no \texttt{Cargo.toml} da raiz. Cada um tem um
papel bem definido:

\begin{itemize}
  \item \texttt{libsurfer}: a biblioteca central. Concentra praticamente toda a
        lógica e a interface. É o crate onde se faz a maior parte das mudanças.
  \item \texttt{surfer}: o binário da aplicação. É pequeno; seu
        \texttt{surfer/src/main.rs} inicializa o \texttt{eframe} e delega para o
        \texttt{libsurfer}. O build deste crate produz o executável (o
        \texttt{surfer.exe} no Windows).
  \item \texttt{surver}: o servidor. Permite abrir os waveforms em uma máquina e
        visualizá-los de outra, sem copiar os arquivos.
  \item \texttt{surfer-translation-types}: os tipos compartilhados entre o núcleo
        e os formatadores/plugins de valor.
  \item \texttt{surfer-wcp}: a implementação do Waveform Control Protocol, o
        canal para controlar o Surfer a partir de fora.
  \item \texttt{surfer-vscode}: a extensão do VS Code (empacota a versão WASM).
  \item \texttt{wasm\_example\_translator}: um exemplo de translator compilado
        para WebAssembly, servindo de modelo para plugins.
\end{itemize}

\subsection{Onde fica cada componente}

Quase tudo vive em \texttt{libsurfer/src}. Os arquivos têm nomes descritivos.
Agrupando por responsabilidade (nomes reais dos arquivos):

\paragraph{Interface (UI).}
\texttt{view.rs} (a view principal), \texttt{drawing\_canvas.rs} (o desenho das
ondas), \texttt{toolbar.rs}, \texttt{statusbar.rs}, \texttt{menus.rs},
\texttt{hierarchy.rs} (a árvore de escopos e sinais), \texttt{command\_prompt.rs}
(a paleta de comandos), \texttt{dialog.rs}, \texttt{overview.rs},
\texttt{tooltips.rs}, \texttt{help.rs}, \texttt{keyboard\_shortcuts.rs},
\texttt{mousegestures.rs}, \texttt{memory\_viewer.rs}, \texttt{logs.rs}.

\paragraph{Estado e lógica.}
\texttt{state.rs} e \texttt{system\_state.rs} (o estado da aplicação),
\texttt{message.rs} (as mensagens/eventos), \texttt{displayed\_item.rs} e
\texttt{displayed\_item\_tree.rs} (os itens exibidos e sua hierarquia),
\texttt{wave\_data.rs} (os dados da waveform aberta), \texttt{viewport.rs},
\texttt{marker.rs}, \texttt{config.rs}. O par \texttt{state\_file\_io.rs} cuida de
salvar e carregar o estado em arquivos \texttt{.surf.ron}, que é justamente o
mecanismo que a AURORA usa para abrir o Surfer já com os sinais curados.

\paragraph{Dados de waveform.}
\texttt{data\_container.rs}, \texttt{transaction\_container.rs} e
\texttt{cxxrtl\_container.rs} (fontes de dados: arquivo, transações, co-simulação
CXXRTL), \texttt{variable\_meta.rs} e demais \texttt{variable\_*.rs}. A leitura
dos formatos é delegada à biblioteca externa \texttt{wellen}.

\paragraph{Desenho e cache.}
\texttt{analog\_renderer.rs} e \texttt{analog\_signal\_cache.rs} (sinais
analógicos), \texttt{clock\_highlighting.rs}, \texttt{graphics.rs},
\texttt{frame\_buffer.rs}, \texttt{arrow.rs}, \texttt{rectangle.rs},
\texttt{item\_drawing\_info.rs}, \texttt{trace\_style.rs}.

\paragraph{Formatadores de valor (translators).}
Ficam em \texttt{libsurfer/src/translation/}. É aqui que um valor de sinal vira
texto legível ou onda. Arquivos: \texttt{basic\_translators.rs} (os tradutores
Bit, Binary, Hexadecimal, entre outros), \texttt{numeric\_translators.rs},
\texttt{fixed\_point.rs}, \texttt{enum\_translator.rs},
\texttt{instruction\_translators.rs}, \texttt{mapping\_translators.rs} (os
mapping translators, usados pela AURORA para decodificar Assembly e linha-fonte),
\texttt{clock.rs}, \texttt{color\_translators.rs}, \texttt{event\_translator.rs},
\texttt{python\_translators.rs} e \texttt{wasm\_translator.rs} (plugins).

\paragraph{Comandos, protocolo e servidor.}
O sistema de comandos está em \texttt{command\_parser.rs},
\texttt{command\_prompt.rs}, \texttt{batch\_commands.rs} e \texttt{fzcmd.rs}. O
WCP tem a subpasta \texttt{libsurfer/src/wcp/} e o crate \texttt{surfer-wcp}. O
modo remoto está em \texttt{libsurfer/src/remote/} e no crate \texttt{surver}.

\paragraph{Assets e recursos.}
Os temas ficam em \texttt{themes/} (arquivos \texttt{.toml} como \texttt{dark+},
\texttt{light+}, \texttt{ibm}, \texttt{okabe-ito}, \texttt{petroff-dark},
\texttt{rose-pine-moon}, além do \texttt{default\_theme.toml} na raiz). Os
waveforms de exemplo ficam em \texttt{examples/}, as imagens de referência dos
testes visuais em \texttt{snapshots/}, e a documentação de usuário em
\texttt{docs/} (subpastas \texttt{commands}, \texttt{configuration},
\texttt{plugins} e \texttt{development}).

\section{O nosso fork: surfer-aurora}

O Surfer-Aurora é um fork do Surfer mantido pela NIPSCERN, em
\url{https://gitlab.com/nips-cern/surfer-aurora}. Ele preserva o vínculo com o
repositório original: dá para puxar as melhorias do upstream e, se quisermos,
propor as nossas de volta. É onde entram as melhorias específicas para uso na
AURORA.

Dois pontos de licença que orientam o fork:

\begin{itemize}
  \item O Surfer é EUPL-1.2, uma licença copyleft. Como a AURORA distribui o
        binário modificado, o fork precisa permanecer público e o código-fonte
        correspondente fica disponível a quem recebe.
  \item A AURORA não é contaminada pela EUPL porque executa o
        \texttt{surfer-aurora.exe} como um processo separado (arm's-length), sem
        linkar o código do Surfer. A AURORA mantém a sua própria licença.
\end{itemize}

O binário do fork se chama \texttt{surfer-aurora.exe}. Esse nome não vem do
Cargo (o crate continua chamado \texttt{surfer}, então o compilador ainda produz
\texttt{surfer.exe}); a renomeação acontece no empacotamento, tanto no build
local quanto no CI. Isso mantém o fonte do fork idêntico ao do upstream e evita
conflitos ao sincronizar.

\section{Como o nosso executável é gerado}

Existem dois caminhos: o build local (para desenvolvimento) e o build no CI (que
publica o binário oficial que a AURORA baixa).

\subsection{Build local}

Requer a toolchain do Rust instalada. No clone do fork:

\begin{lstlisting}[language=bash]
git clone https://gitlab.com/nips-cern/surfer-aurora.git
cd surfer-aurora
cargo build --release
\end{lstlisting}

Isso produz \texttt{target/release/surfer.exe}. Para usar na AURORA, copie e
renomeie:

\begin{lstlisting}[language=bash]
cp target/release/surfer.exe \
   ../aurora/components/Packages/surfer/surfer-aurora.exe
\end{lstlisting}

\subsection{Build no CI (o que publica o binário oficial)}

O binário oficial é gerado por um pipeline dedicado do fork, definido no arquivo
\texttt{.gitlab-ci-aurora.yml}. Esse arquivo é separado do
\texttt{.gitlab-ci.yml} do upstream de propósito: ele está apontado como o
``CI/CD configuration file'' do projeto (em Settings, CI/CD), então sincronizar o
upstream nunca gera conflito nele, e não rodamos os jobs de macOS/ARM/docs do
upstream, para os quais não temos runners.

O que o pipeline faz, disparado apenas no push de uma tag:

\begin{enumerate}
  \item Roda em um runner Linux (imagem \texttt{rust:latest}) e faz
        cross-compile para Windows. Não usa runner Windows: instala o
        \texttt{mingw-w64} e compila para o alvo \texttt{x86\_64-pc-windows-gnu}.
        Essa é a mesma receita do build de release do upstream.
  \item O comando central é, em essência:
\begin{lstlisting}[language=bash]
rustup target add x86_64-pc-windows-gnu
cargo build --target x86_64-pc-windows-gnu --release --locked --features accesskit
\end{lstlisting}
  \item Renomeia o binário para \texttt{surfer-aurora.exe} e o empacota em
        \texttt{surfer-aurora\_win\_<tag>.zip}.
  \item Calcula o SHA-256 do zip (registrado no log e no artefato
        \texttt{SHA256.txt}).
  \item Publica o zip no registro de pacotes genérico do próprio fork, via
        \texttt{JOB-TOKEN}. Como o projeto é público, esse pacote pode ser
        baixado sem autenticação.
  \item Cria um Release do GitLab na tag, com um link para o pacote.
\end{enumerate}

\subsection{Passo a passo para gerar uma nova versão}

Duas fases. Primeiro, disparar o build no fork:

\begin{lstlisting}[language=bash]
cd surfer-aurora
git tag v0.7.0-nips.3          # o esquema de versao e nosso
git push origin v0.7.0-nips.3  # o push da tag dispara o pipeline
\end{lstlisting}

O CI compila e publica (leva por volta de 20 minutos a frio) e gera o
\texttt{SHA256.txt} no job \texttt{build-windows}. Segundo, apontar a AURORA para
o novo pacote: no arquivo \texttt{components/Scripts/download-surfer.js}, atualize
o objeto \texttt{FORK\_ARTIFACT} (a tag, o nome do zip, a URL e o \texttt{sha256}
lido do \texttt{SHA256.txt}). A partir daí, \texttt{npm run bootstrap} baixa o
binário novo sozinho.

O esquema de versão que adotamos preserva a base do upstream: \texttt{v0.7.0}
(a base) mais uma revisão da NIPSCERN, por exemplo \texttt{v0.7.0-nips.2},
incrementando a cada release nosso e re-baseando quando o upstream subir.

\subsection{Versão publicada atualmente}

\begin{itemize}
  \item Tag: \texttt{v0.7.0-nips.2}
  \item Pacote: \texttt{surfer-aurora\_win\_v0.7.0-nips.2.zip} (contém
        \texttt{surfer-aurora.exe})
  \item SHA-256:
        \texttt{a27dac113e0b74b2bcc6ec477beab040102ad1889720b7d8b15047cd37efa8ba}
  \item Projeto (id numérico usado na URL do registro): \texttt{84576006}
\end{itemize}

\section{Fluxo de trabalho: do clone ao novo binário}

Esta seção é um passo a passo concreto do ciclo completo de desenvolvimento no
fork: clonar, modificar, testar, commitar e gerar o novo executável. Requer a
toolchain do Rust e o \texttt{git} instalados.

\subsection{Clonar e preparar os remotes}

Clone o fork e adicione o repositório original como \texttt{upstream}, para poder
puxar as melhorias deles depois. O push para o \texttt{upstream} é desabilitado
de propósito: o caminho de volta é por Merge Request, não por push direto.

\begin{lstlisting}[language=bash]
git clone https://gitlab.com/nips-cern/surfer-aurora.git
cd surfer-aurora
git remote add upstream https://gitlab.com/surfer-project/surfer.git
git remote set-url --push upstream DISABLED
git remote -v   # confere: origin = fork, upstream = surfer-project
\end{lstlisting}

Para o \texttt{git push} funcionar sem abrir prompt, a máquina precisa de
credencial do GitLab. Uma opção é usar a CLI \texttt{glab} como credential
helper; outra é configurar um token de acesso pessoal. Basta uma vez por
máquina.

\subsection{Criar um branch e modificar}

Trabalhe sempre em um branch dedicado, não direto na \texttt{main}. Isso mantém a
mudança isolada e facilita tanto o merge quanto uma eventual proposta de volta
para o upstream.

\begin{lstlisting}[language=bash]
git checkout -b nips-minha-melhoria
\end{lstlisting}

Quase tudo que se altera fica em \texttt{libsurfer/src/} (ver Seção 3.3). Prefira
mudanças pequenas e localizadas: quanto menor o diff em relação ao upstream,
menor o atrito ao sincronizar depois.

\subsection{Testar}

O ciclo de teste é todo pelo Cargo. Do mais rápido ao mais completo:

\begin{lstlisting}[language=bash]
cargo check          # compila e checa o codigo, sem gerar binario (rapido)
cargo run --release -- examples/analog.vcd   # abre a GUI num waveform de exemplo
cargo test           # roda a suite de testes (inclui testes de snapshot)
cargo clippy         # lints
cargo fmt            # formatacao (rustfmt.toml na raiz)
\end{lstlisting}

O \texttt{cargo run} abre o Surfer de verdade, então dá para ver a mudança na
tela. A pasta \texttt{examples/} tem vários arquivos de onda para testar (por
exemplo \texttt{examples/analog.vcd}). Os testes de snapshot comparam imagens de
referência em \texttt{snapshots/}; se você mudou algo visual de propósito, há
scripts na raiz (\texttt{accept\_snapshots.bash}) para atualizar as referências.

\subsection{Commitar e enviar}

Com a mudança testada, commite no branch e envie para o fork:

\begin{lstlisting}[language=bash]
git add -A
git commit -m "feat: descricao curta da melhoria"
git push origin nips-minha-melhoria
\end{lstlisting}

No GitLab, abra um Merge Request do branch para a \texttt{main} do fork (ou faça
o merge localmente). Se a melhoria também interessa ao projeto original, dá para
abrir um Merge Request do fork para o \texttt{upstream} a partir daí.

\subsection{Gerar o novo .exe}

Para um teste rápido na sua máquina, o build local basta (Seção 5.1): gera
\texttt{target/release/surfer.exe}, que você renomeia para
\texttt{surfer-aurora.exe} e coloca em \texttt{components/Packages/surfer/} da
AURORA.

Para o binário oficial que a AURORA baixa, é preciso que a mudança esteja na
\texttt{main} e então cortar uma tag, que dispara o pipeline (Seção 5.2):

\begin{lstlisting}[language=bash]
git checkout main
git merge nips-minha-melhoria       # traz a melhoria para a main
git push origin main
git tag v0.7.0-nips.3               # nova versao
git push origin v0.7.0-nips.3       # o push da tag dispara o CI
\end{lstlisting}

O CI compila e publica o \texttt{surfer-aurora\_win\_v0.7.0-nips.3.zip} no
registro de pacotes, com o \texttt{SHA256.txt}. Por fim, aponte a AURORA para a
nova versão editando o \texttt{FORK\_ARTIFACT} em
\texttt{components/Scripts/download-surfer.js} (tag, nome do zip, URL e o
\texttt{sha256}). Depois disso, \texttt{npm run bootstrap} instala o binário novo
automaticamente.

\section{Como a AURORA usa o surfer-aurora.exe}

\subsection{Download automático}

O script \texttt{components/Scripts/download-surfer.js} baixa o zip do registro
de pacotes do fork, confere o SHA-256, extrai o \texttt{surfer-aurora.exe} e o
coloca em \texttt{components/Packages/surfer/}. Como o projeto é público, o
download é feito sem autenticação, igual aos outros toolchains. Esse script roda
como parte do \texttt{npm run bootstrap}. O binário em si não é versionado no
repositório da AURORA; vem sempre do build/CI.

\subsection{Onde a AURORA resolve o binário}

A AURORA procura exatamente por \texttt{surfer-aurora.exe} em dois pontos:

\begin{itemize}
  \item \texttt{js/compilation/wave\_toolchain.js} resolve o caminho
        \texttt{components/Packages/surfer/surfer-aurora.exe}.
  \item \texttt{main/compile/binary\_allowlist.js} autoriza esse nome de binário
        (gate de segurança de execução).
\end{itemize}

\subsection{Como é lançado}

O viewer é lançado como um processo externo (arm's-length), pela chamada
\texttt{launch-surfer} no processo principal (\texttt{main/ipc/compile.js}). A
AURORA passa o arquivo de onda e, opcionalmente, um arquivo de layout que faz o
Surfer abrir já com os sinais curados (cores, formatos, grupos, marcadores). Esse
layout é um \texttt{.surf.ron} gerado por
\texttt{js/wave/surfer\_layout\_writer.ts}, e é entregue ao Surfer pela opção
\texttt{-s}. Em linha de comando, é equivalente a:

\begin{lstlisting}[language=bash]
surfer-aurora.exe caminho/para/onda.fst -s caminho/para/layout.surf.ron
\end{lstlisting}

Se o binário não estiver presente, a AURORA cai de forma limpa para o GTKWave,
então o botão de onda sempre produz um viewer.

\subsection{Uso standalone}

O \texttt{surfer-aurora.exe} também é um Surfer completo e pode ser aberto por
fora da AURORA, apontando direto para um arquivo de onda:

\begin{lstlisting}[language=bash]
surfer-aurora.exe onda.vcd
surfer-aurora.exe onda.fst
\end{lstlisting}

\section{Sincronizar com o upstream}

O fluxo para trazer as melhorias do projeto original está montado no clone, com
dois remotes: \texttt{origin} apontando para o nosso fork e \texttt{upstream}
para o repositório do surfer-project (o push para o upstream fica desabilitado,
já que o caminho de volta é por Merge Request). Puxar as melhorias deles:

\begin{lstlisting}[language=bash]
git fetch upstream
git merge upstream/main
git push origin main
\end{lstlisting}

Como a AURORA usa um binário compilado, puxar o código não muda o executável
sozinho: depois do merge é preciso gerar uma nova versão (Seção 5.3) para que a
melhoria chegue ao \texttt{surfer-aurora.exe}.

\section{Licença e créditos}

O Surfer é desenvolvido pelo surfer-project sob a licença EUPL-1.2. O
Surfer-Aurora é um fork derivado desse trabalho, mantido público sob a mesma
licença. A atribuição ao projeto original e ao nosso fork consta nos arquivos
\texttt{LICENSE} e \texttt{THIRD\_PARTY\_NOTICES.md} da AURORA, junto do texto da
EUPL-1.2 preservado no fork.

\end{document}
